\section{Rを使ったt検定と検定力分析}
\label{b16_ttest_R}

\subsection{授業内容}
\subsubsection{概要}
夏休みを挟んで，前期に学んだ帰無仮説検定の基本的な考え方を復習し，統計ソフトウェアRを用いた実践的な分析方法を学ぶ。特に，心理学研究でよく用いられるt検定(対応のあるt検定と対応のないt検定)に焦点を当て，Rによる実行方法とその結果の解釈について理解を深める。さらに，シミュレーションを通じて効果量とサンプルサイズが検定結果に与える影響を体験的に学び，適切な研究デザインの重要性について理解する。また，検定力分析の基本的な考え方と実行方法についても学ぶ。

\subsubsection{コマ主題細目}
\begin{description}
	\item[帰無仮説検定の復習] 前期に学んだ帰無仮説検定の基本的な考え方と手順について復習する。帰無仮説と対立仮説の設定，検定統計量の算出，p値の解釈，有意水準に基づく判断という一連の流れを確認する。特に，検定結果の解釈と報告方法について理解を深め，統計的有意性と実質的重要性の違いについても再確認する。また，t分布の特徴と，標準正規分布との関係についても理解する。
	      \begin{flushright}
		      $\to$ \textcite{Kawabata201408}のP.81--96，\textcite{Yamada2004}のP.104--119.
	      \end{flushright}

	\item[対応のあるt検定] 対応のあるt検定(一要因被験者内計画の2水準間の比較)の考え方と実行方法について学ぶ。この検定は同一被験者から得られた2つの条件下でのデータ(事前・事後測定など)を比較する場合に用いられる。Rでの実行方法として，`t.test()`関数の使い方，特に`paired=TRUE`オプションの意味と使用方法を理解する。また，結果の読み方と解釈，効果量の算出方法についても学ぶ。実際のデータセットを使用した演習を通じて，対応のあるt検定の実践的な実行方法を身につける。
	      \begin{flushright}
		      $\to$ \textcite{kosugi2019}のP.126--132，\textcite{Shimizu2021}のP.124--130.
	      \end{flushright}

	\item[対応のないt検定] 対応のないt検定(独立した2群の平均値の比較)の考え方と実行方法について学ぶ。この検定は異なる被験者群から得られたデータを比較する場合に用いられる。Rでの実行方法として，`t.test()`関数の基本的な使い方と，分散の等質性を仮定する場合(`var.equal=TRUE`)と仮定しない場合(Welchの補正)の違いについて理解する。また，検定結果の読み方と解釈，効果量の算出方法についても学ぶ。実際のデータセットを使用した演習を通じて，対応のないt検定の実践的な実行方法を身につける。
	      \begin{flushright}
		      $\to$ \textcite{kosugi2019}のP.132--138，\textcite{Shimizu2021}のP.130--136.
	      \end{flushright}

	% \item[効果量とサンプルサイズ] Rを用いたシミュレーションを通じて，効果量とサンプルサイズが検定結果に与える影響について体験的に学ぶ。異なる効果量(小・中・大)とサンプルサイズの組み合わせでデータを生成し，t検定を実行することで，どのような条件で統計的有意性が得られやすいか，または得られにくいかを確認する。これにより，効果量が小さい場合は大きなサンプルサイズが必要であること，サンプルサイズが小さい場合は検出力が低下することなどを理解する。また，シミュレーションの結果から，標本サイズを適切に設定することの重要性についても学ぶ。
	%       \begin{flushright}
	% 	      $\to$ 効果量については\textcite{Matia201201}のP.43--46，シミュレーションの基本的な考え方については\textcite{Kinosady2021}のP.178--186.
	%       \end{flushright}

	\item[検定力分析の基礎] 検定力(第2種の誤りを犯さない確率)の概念と，研究デザインにおける重要性について学ぶ。特に，Rの`power.t.test()`関数を用いた検定力分析の実行方法と解釈について理解する。事前の検定力分析(研究計画段階での必要サンプルサイズの算出)と事後の検定力分析(研究結果の評価)の違いと用途について学び，研究デザインの質を高めるための手法として検定力分析を活用する方法を身につける。また，サンプルサイズ，効果量，有意水準，検定力の関係について理解を深める。
	      \begin{flushright}
		      $\to$ 検定力分析については\textcite{Kawabata201408}のP.90--93，\textcite{kosugi2023}のP.187--224も参照。
	      \end{flushright}

\end{description}
\subsubsection{キーワード}
\begin{itemize}
	\item 対応のあるt検定と対応のないt検定
	\item Rによる統計分析の実行
	\item 効果量とCohen's d
	\item サンプルサイズとシミュレーション
	\item 検定力分析
\end{itemize}
\subsection{授業情報}
\paragraph{コマの展開方法}
講義と演習
\paragraph{標準シラバスにおける位置づけ}
科目番号5；心理学統計法；1.心理学で用いられる統計手法；G 推測統計：t検定
\subsubsection{予習・復習課題}
\paragraph{予習}
前期に学んだ帰無仮説検定の基本的な考え方と手順について復習しておくこと。特に，帰無仮説と対立仮説の設定，検定統計量の算出，p値の解釈，有意水準に基づく判断という一連の流れを確認しておくとよい。また，RとRStudioの基本的な操作方法についても再確認しておくこと。前期に作成したスクリプトやRmarkdownファイルを見直し，基本的なコマンドや関数の使い方を思い出しておくことが望ましい。

\paragraph{復習}
授業で学んだ内容を定着させるために，以下の課題に取り組むこと：

\begin{enumerate}
\item 提供されたデータセットを用いて，対応のあるt検定と対応のないt検定をそれぞれ実行し，結果をRmarkdownで適切にレポートにまとめる
\item 異なる効果量とサンプルサイズの組み合わせでシミュレーションデータを生成し，t検定を実行して結果を比較する
\item 与えられた研究シナリオに対して，適切なサンプルサイズを算出するための検定力分析を実行する
\item 検定結果の統計的有意性と効果量の大きさの両方を考慮した解釈を行い，研究における意義について考察する
\end{enumerate}

これらの課題をRmarkdownにまとめ，提出すること。また，検定力分析とサンプルサイズの重要性について理解を深めるために，\textcite{Matia201201}や\textcite{Kawabata201408}，\textcite{kosugi2023}の関連する章を読むことが推奨される。